Figure~\ref{fig:atlas-supremum} contrasts a supremum that is not attained with a maximum that belongs to the set.

\begin{figure}[!htbp]
\centering
\begin{tikzpicture}[>=Stealth]
\draw[->,black!45] (-0.3,0) -- (6.5,0);
\draw[figblue,very thick] (0,0) -- (5,0);
\fill[figblue] (0,0) circle (2pt);
\draw[figblue,thick,fill=white] (5,0) circle (2.5pt);
\node[below] at (0,0) {$0$};
\node[below] at (5,0) {$1$};
\node[left] at (-0.5,0) {$[0,1)$};
\draw[dashed,figgreen] (3.5,-0.2) -- (3.5,0.7);
\node[above,figgreen] at (3.5,0.7) {$1-\varepsilon$};
\fill[figred] (4.25,0) circle (2.5pt);
\draw[figred] (4.25,0.1) -- (4.25,1.3);
\node[above,figred] at (4.25,1.3) {$1-\varepsilon/2$};
\draw[->,black!45] (-0.3,-1.7) -- (6.5,-1.7);
\draw[figblue,very thick] (0,-1.7) -- (5,-1.7);
\fill[figblue] (0,-1.7) circle (2pt) (5,-1.7) circle (2.5pt);
\node[below] at (0,-1.7) {$0$};
\node[below] at (5,-1.7) {$1$};
\node[left] at (-0.5,-1.7) {$[0,1]$};
\end{tikzpicture}
\caption{The sets $[0,1)$ (top) and $[0,1]$ (bottom) both have supremum $1$. Only the second contains that endpoint and therefore has a maximum. For the illustrated choice $0<\varepsilon<1$, the highlighted point $1-\varepsilon/2$ lies within $\varepsilon$ of the supremum in the first set.}
\label{fig:atlas-supremum}
\end{figure}